\section{Exterior Powers and Exterior Algebra}\label{sec:exterior-algebra}
Finite tensor products represent multilinear maps.  Exterior powers impose
one further relation: a multilinear expression should vanish whenever two
inputs agree.  This packages alternating forms, oriented volume, and the
determinant into one construction.

Throughout this section \(R\) is a commutative ring and \(M\) is an
\(R\)-module.  A multilinear map
\[
f:M^k\longrightarrow N
\]
is \emph{alternating} if
\[
f(m_1,\ldots,m_k)=0
\]
whenever two arguments coincide.  This is the same alternating condition
used for determinants in \hyperref[sec:determinants]{Section~6.12}; it is
not to be confused with skew-symmetry in characteristic \(2\).

\begin{theorem}[Exterior-power universal property]\label{thm:exterior-power}
For every integer \(k\geq1\), there is an \(R\)-module
\[
\bigwedge_R^kM
\]
and an alternating multilinear map
\[
\omega_k:M^k\longrightarrow\bigwedge_R^kM,
\qquad
(m_1,\ldots,m_k)\longmapsto m_1\wedge\cdots\wedge m_k,
\]
such that every alternating multilinear map
\[
f:M^k\longrightarrow N
\]
has a unique linear factorization
\[
\widetilde f:\bigwedge_R^kM\longrightarrow N.
\]
\end{theorem}
\[
\begin{tikzcd}[column sep=large, row sep=large]
M^k \arrow[r, "\omega_k"] \arrow[dr, "f"'] &
\bigwedge_R^k M \arrow[d, "\widetilde f"] \\
& N
\end{tikzcd}
\]
\begin{proof}
By Theorem~\ref{thm:finite-tensors}, form the finite tensor power
\[
M^{\otimes k}
=M\otimes_R\cdots\otimes_RM.
\]
Let \(J\) be the submodule generated by all pure tensors having two equal
entries, and define
\[
\bigwedge_R^kM=M^{\otimes k}/J.
\]
The quotient of the universal multilinear tensor map is \(\omega_k\), and
it is alternating by construction.  If \(f\) is alternating, its associated
linear map
\[
M^{\otimes k}\longrightarrow N
\]
vanishes on every generator of \(J\), so it descends uniquely through the
quotient.  This gives the required factorization.  Any two modules with this
property are canonically isomorphic by applying their universal properties
to one another's universal alternating maps.
\end{proof}

\begin{proposition}[Alternating wedge identities]\label{prop:wedge-signs}
If two entries agree, then
\[
m_1\wedge\cdots\wedge m_k=0.
\]
Interchanging two entries changes sign; consequently
\[
m_{\sigma(1)}\wedge\cdots\wedge m_{\sigma(k)}
=\operatorname{sgn}(\sigma)\,
m_1\wedge\cdots\wedge m_k
\]
for every \(\sigma\in S_k\).
\end{proposition}
\begin{proof}
The first assertion is part of the construction.  Holding all other entries
fixed, alternation and multilinearity give
\[
0=\omega_k(\ldots,x+y,\ldots,x+y,\ldots).
\]
The repeated-\(x\) and repeated-\(y\) terms vanish, leaving the sum of the
two wedges with \(x\) and \(y\) interchanged.  Thus a transposition
contributes \(-1\).  The sign theory of Chapter~2 then gives the formula for
an arbitrary permutation.
\end{proof}

\noindent\textbf{Example.} If \(M=R^3\) has basis \(e_1,e_2,e_3\), then
\[
\bigwedge_R^2M
\]
has basis
\[
e_1\wedge e_2,\qquad e_1\wedge e_3,\qquad e_2\wedge e_3.
\]
Terms \(e_i\wedge e_i\) vanish, and reversing two basis vectors changes the
sign, so these are precisely the increasing-index wedges.

\begin{theorem}[Basis of an exterior power]\label{thm:exterior-basis}
If \(M\) is finite free with ordered basis
\[
(e_1,\ldots,e_n),
\]
then the wedges
\[
e_{i_1}\wedge\cdots\wedge e_{i_k},
\qquad
1\leq i_1<\cdots<i_k\leq n,
\]
form a basis of
\[
\bigwedge_R^kM.
\]
Hence
\[
\operatorname{rank}_R\bigl(\bigwedge_R^kM\bigr)=\binom nk.
\]
\end{theorem}
\begin{proof}
Multilinearity and Proposition~\ref{prop:wedge-signs} show that these wedges
span: expand every argument in the basis, discard repeated-index terms, and
reorder the rest.  To prove independence, fix an increasing \(k\)-tuple
\[
I=(i_1,\ldots,i_k).
\]
For \(m_j=\sum_ra_{rj}e_r\), define
\[
f_I(m_1,\ldots,m_k)=
\det(a_{i_rj})_{1\leq r,j\leq k}.
\]
The determinant theory of
\hyperref[sec:determinants]{Section~6.12} shows that \(f_I\) is alternating
and multilinear, so it induces a linear map from \(\bigwedge_R^kM\) to
\(R\).  On the displayed wedges it has value \(1\) for the tuple \(I\) and
\(0\) for every other increasing tuple.  Applying these maps to a linear
relation proves that every coefficient is zero.
\end{proof}

We set
\[
\bigwedge_R^0M=R,
\qquad
\bigwedge_R^1M=M.
\]
The basis theorem gives
\[
\bigwedge_R^kM=0
\qquad(k>n)
\]
when \(M\) is finite free of rank \(n\), and it shows that
\[
\bigwedge_R^nM
\]
is free of rank one.

\begin{proposition}[Induced exterior maps]\label{prop:exterior-maps}
An \(R\)-linear map
\[
T:M\longrightarrow N
\]
induces a unique linear map
\[
\bigwedge^kT:\bigwedge_R^kM\longrightarrow\bigwedge_R^kN
\]
satisfying
\[
\bigwedge^kT(m_1\wedge\cdots\wedge m_k)
=Tm_1\wedge\cdots\wedge Tm_k.
\]
Moreover,
\[
\bigwedge^k(\operatorname{id}_M)
=\operatorname{id}_{\bigwedge_R^kM},
\qquad
\bigwedge^k(S\circ T)
=\bigl(\bigwedge^kS\bigr)\circ\bigl(\bigwedge^kT\bigr).
\]
\end{proposition}
\begin{proof}
The map
\[
(m_1,\ldots,m_k)\longmapsto Tm_1\wedge\cdots\wedge Tm_k
\]
is alternating and multilinear.  By Theorem~\ref{thm:exterior-power}, it
therefore induces a unique map.  The two displayed identities agree
on every pure wedge, and pure wedges generate the relevant exterior powers.
\end{proof}

\begin{theorem}[Top exterior action]
Let \(V\) be \(n\)-dimensional over \(F\), and let \(T:V\to V\).  Then
\[
\bigwedge^nT
\]
acts on the one-dimensional space
\[
\bigwedge_F^nV
\]
by multiplication by \(\det(T)\).  Equivalently,
\[
Tv_1\wedge\cdots\wedge Tv_n
=\det(T)\,
\bigl(v_1\wedge\cdots\wedge v_n\bigr).
\]
\end{theorem}
\begin{proof}
Fix an ordered basis \((v_1,\ldots,v_n)\).  The linear functional on
\(\bigwedge_F^nV\) that takes \(v_1\wedge\cdots\wedge v_n\) to \(1\)
corresponds, by the exterior universal property, to a normalized alternating
\(n\)-linear form on \(V\).  By the uniqueness theorem for determinants in
\hyperref[sec:determinants]{Section~6.12}, its value on
\[
(Tv_1,\ldots,Tv_n)
\]
is \(\det(T)\).  Since the top exterior power is one-dimensional, this is
exactly the scalar by which \(\bigwedge^nT\) acts.
\end{proof}

This gives a later, equivalent characterization of the determinant already
developed from normalized alternating multilinear forms in
\hyperref[sec:determinants]{Section~6.12}.  Since
\[
\dim_F\bigwedge_F^nV=1,
\]
the determinant of \(T\) is the unique scalar \(\lambda\in F\) such that
\[
\bigwedge^nT
=\lambda\operatorname{id}_{\bigwedge_F^nV}.
\]
Equivalently, for every ordered basis \((v_1,\ldots,v_n)\) of \(V\),
\[
Tv_1\wedge\cdots\wedge Tv_n
=\det(T)\bigl(v_1\wedge\cdots\wedge v_n\bigr).
\]
Thus this top-exterior-power viewpoint describes the same invariant as the
Leibniz and alternating-multilinear descriptions, rather than defining a
new determinant.

For \(A\in M_n(F)\), regarded as an endomorphism of \(F^n\), the matrix of
\[
\bigwedge^nA
\]
with respect to the one-element basis
\[
e_1\wedge\cdots\wedge e_n
\]
of \(\bigwedge_F^nF^n\) is the \(1\times1\) matrix
\[
(\det A).
\]
Moreover, the composition law proved in Proposition~\ref{prop:exterior-maps}
makes determinant multiplicativity conceptually immediate:
\[
\bigwedge^n(S\circ T)
=\bigl(\bigwedge^nS\bigr)\circ\bigl(\bigwedge^nT\bigr)
\]
implies
\[
\det(S\circ T)=\det(S)\det(T).
\]
This is a retrospective explanation of the multiplicativity already proved
in Chapter~6, not a replacement for that earlier proof.

\section{The Exterior Algebra and Alternating Forms}\label{sec:exterior-forms}
The individual exterior powers fit together into a graded algebra.  This
uses an arbitrary direct sum of already constructed modules; it does
\emph{not} introduce an infinite tensor product.

\begin{definition}
The \emph{exterior algebra} of \(M\) is
\[
\bigwedge M
=\bigoplus_{k\geq0}\bigwedge_R^kM.
\]
For pure wedges, define the product by concatenation:
\[
\begin{aligned}
&(m_1\wedge\cdots\wedge m_p)
\wedge
(n_1\wedge\cdots\wedge n_q)\\
&\qquad =
m_1\wedge\cdots\wedge m_p
\wedge n_1\wedge\cdots\wedge n_q.
\end{aligned}
\]
\end{definition}

The right-hand side is alternating separately in the \(m\)'s and in the
\(n\)'s, so the two exterior universal properties make this a well-defined
bilinear map
\[
\bigwedge_R^pM\times\bigwedge_R^qM
\longrightarrow
\bigwedge_R^{p+q}M.
\]
Pure wedges generate each factor, so associativity follows from associativity
of concatenation on pure wedges.  The degree-\(k\) summand is the \(k\)-th
graded piece.

\begin{proposition}[Graded commutativity]
For homogeneous elements
\[
\alpha\in\bigwedge_R^pM,
\qquad
\beta\in\bigwedge_R^qM,
\]
one has
\[
\alpha\wedge\beta=(-1)^{pq}\beta\wedge\alpha.
\]
\end{proposition}
\begin{proof}
For pure wedges, moving the block of \(p\) entries past the block of \(q\)
entries requires \(pq\) transpositions.  Proposition~\ref{prop:wedge-signs}
gives the sign.  Bilinearity extends the identity to arbitrary homogeneous
elements.  The argument uses alternation itself, so it remains valid in
characteristic \(2\), where the displayed sign may equal \(1\).
\end{proof}

\begin{theorem}[Alternating forms as exterior powers of the dual]
Let \(V\) be a finite-dimensional vector space over \(F\).  There is a
natural isomorphism from
\[
\bigwedge_F^kV^\vee
\]
onto the space of alternating \(k\)-linear maps
\[
V^k\longrightarrow F.
\]
On a decomposable element it is given by
\[
\varphi_1\wedge\cdots\wedge\varphi_k
\longmapsto
\bigl((v_1,\ldots,v_k)\mapsto
\det(\varphi_i(v_j))_{1\leq i,j\leq k}\bigr).
\]
\end{theorem}
\begin{proof}
The displayed determinant is alternating and multilinear in the vectors, so
the exterior universal property makes the assignment well-defined and
linear.  If \((v_1,\ldots,v_n)\) is a basis of \(V\), then the wedges of its
dual basis indexed by increasing \(k\)-tuples map to the corresponding
coordinate alternating forms.  These forms are linearly independent and
span by evaluating an alternating form on basis tuples.  Thus the map is an
isomorphism.
\end{proof}

This identifies the algebra developed here with the pointwise algebra behind
differential forms.  On a smooth manifold, a differential \(k\)-form assigns
to each point \(p\) an alternating \(k\)-linear form on the tangent space
\[
T_pM,
\]
that is, an element of
\[
\bigwedge^k(T_pM)^\vee.
\]
The wedge product sends a \(p\)-form and a \(q\)-form at \(p\) to an element
of
\[
\bigwedge^{p+q}(T_pM)^\vee.
\]
No manifold theory is needed here; the point is that exterior algebra is the
linear-algebraic model that makes those geometric objects possible.
